\section*{PRÁCTICA No. 6}

\subsection*{APLICACIÓN DE LA LEY DEL CUADRADO INVERSO}

\textbf{OBJETIVO:} El alumno debe saber aplicar la Ley del Cuadrado Inverso a la rapidez de exposición esperada, a partir de una fuente radiactiva colocada en un lugar conocido.

\textbf{MATERIAL:}
\begin{itemize}
    \item Equipo portátil Detector de Radiación tipo Geiger-Müller.
    \item Contenedor de Gammagrafía. (Conteniendo una fuente radiactiva con la suficiente actividad como para poder tener en sus costados niveles de rapidez de exposición de alrededor de 50 a 100 Mr/h).
\end{itemize}

% Commented section: The diagram depicted in the image showing a marked intensity  versus distance graph which cannot be accurately transcribed into LaTeX without the actual data points.
% Tip: It would be beneficial to include a recreation of this graph using a LaTeX package like pgfplots, based on actual measurements or approximations from practice details.